\documentclass[11pt]{article}
\usepackage[T1]{fontenc}
\usepackage[utf8]{inputenc}
\usepackage{lmodern}
\usepackage{amsmath,amssymb,mathtools,array,booktabs}
\usepackage{geometry}
\usepackage{microtype}
\geometry{a4paper,margin=1.45cm}
\setlength{\parindent}{1.25em}
\setlength{\parskip}{0pt}
\allowdisplaybreaks
\setlength{\emergencystretch}{30em}
\sloppy
\hbadness=10000
\vbadness=10000
\hfuzz=10000pt
\vfuzz=10000pt
\raggedbottom

\providecommand{\mG}{\mathfrak{G}}
\providecommand{\mA}{\mathfrak{A}}
\providecommand{\mB}{\mathfrak{B}}
\providecommand{\mC}{\mathfrak{C}}
\providecommand{\mD}{\mathfrak{D}}
\providecommand{\mE}{\mathfrak{E}}
\providecommand{\mH}{\mathfrak{H}}
\providecommand{\mK}{\mathfrak{K}}
\providecommand{\mM}{\mathfrak{M}}
\providecommand{\mN}{\mathfrak{N}}
\providecommand{\mR}{\mathfrak{R}}
\providecommand{\mS}{\mathfrak{S}}
\providecommand{\mT}{\mathfrak{T}}
\providecommand{\mU}{\mathfrak{U}}
\providecommand{\mO}{\mathfrak{O}}
\providecommand{\mo}{\mathfrak{o}}
\providecommand{\ma}{\mathfrak{a}}
\providecommand{\mb}{\mathfrak{b}}
\providecommand{\mc}{\mathfrak{c}}
\providecommand{\md}{\mathfrak{d}}
\providecommand{\mf}{\mathfrak{f}}
\providecommand{\mh}{\mathfrak{h}}
\providecommand{\ml}{\mathfrak{l}}
\providecommand{\mn}{\mathfrak{n}}
\providecommand{\mr}{\mathfrak{r}}
\providecommand{\mt}{\mathfrak{t}}
\providecommand{\mx}{\mathfrak{x}}
\providecommand{\mz}{\mathfrak{z}}
\providecommand{\mZ}{\mathfrak{Z}}
\providecommand{\mL}{\mathfrak{L}}
\providecommand{\rank}{\operatorname{rank}}
\providecommand{\iso}{\simeq}
\providecommand{\ann}{\operatorname{Ann}}
\providecommand{\mat}[1]{\begin{pmatrix}#1\end{pmatrix}}
\usepackage[ngerman]{babel}
\begin{document}


\section*{34. Hyperkomplexe Größen und Darstellungstheorie}
\emph{Fortsetzung und Abschluss: III. Kapitel, §§15--19, und IV. Kapitel, §§20--26.}

\subsection*{III. Kapitel. Modul- und Darstellungstheorie}

\subsection*{§ 15. Darstellungen und Darstellungsmoduln}
Sei $\mo$ ein Ring, $K$ ein Ring mit Einheitselement. (Bei späteren Anwendungen ist $K$ immer ein Körper.)

Eine Darstellung $n$-ten Grades von $\mo$ in $K$ ist eine Homomorphie
\[
        \mo\sim\mathfrak D,
\]
wo $\mathfrak D$ ein Ring aus Matrizes $n$-ten Grades in $K$ ist.

Unter einem Darstellungsmodul von $\mo$ in bezug auf $K$ versteht man einen Doppelmodul $\mM$ (§1, 5), der Links-$\mo$-Modul und Rechts-$K$-Modul ist:
\[
        \mo\mM\subseteq\mM,
        \qquad
        \mM K\subseteq\mM,
\]
der weiter direkte Summe von endlich vielen eingliedrigen $K$-Moduln ist:
\[
        \mM=x_1K+\cdots+x_nK;
\]
und wo das Einheitselement von $K$ Einheitsoperator ist.

Jeder Darstellungsmodul führt zu einer Darstellung. Sei $c$ in $\mo$ und
\[
        cx_k=\sum_i x_i\gamma_{ik},
\]
oder
\[
        c(x_1,\ldots,x_n)=(x_1,\ldots,x_n)C,
        \qquad C=(\gamma_{ik}).
\]
Dann bilden die Matrizes $C$ eine Darstellung für $c$, denn
\[
        (b+c)x_k=\sum_i x_i(\beta_{ik}+\gamma_{ik}),
\]
\[
\begin{aligned}
        bcx_k
        &=b\sum_jx_j\gamma_{jk}
          =\sum_j b x_j\gamma_{jk}
          =\sum_j\sum_i x_i\beta_{ij}\gamma_{jk}  \\
        &=\sum_i x_i\biggl(\sum_j\beta_{ij}\gamma_{jk}\biggr),
\end{aligned}
\]
oder
\[
        bc(x_1,\ldots,x_n)=(x_1,\ldots,x_n)BC.
\]

Umgekehrt: Jede Darstellung $\mathfrak D$ von $\mo$ gehört zu einem Darstellungsmodul, und zwar zu einer bestimmten Basis desselben. Man verstehe unter $\mM$ nämlich die Gesamtheit der formalen Linearformen in $x_1,\ldots,x_n$,
\[
        y=\sum_i x_i\alpha_i.
\]
Dann ist $\mM$ ein Rechts-$K$-Modul. Man definiere weiter, wenn dem Element $c$ die Matrix $C=(\gamma_{ik})$ zugeordnet ist,
\[
        cx_k=\sum_i x_i\gamma_{ik},
        \qquad
        c\biggl(\sum_k x_k\alpha_k\biggr)=\sum_i x_i\sum_k\gamma_{ik}\alpha_k.
\]
Aus den Homomorphierelationen
\[
        c+d\sim C+D,
        \qquad cd\sim CD
\]
folgt
\[
        (c+d)x_i=cx_i+dx_i,
        \qquad cdx_i=c(dx_i),
\]
und dasselbe für Summen $\sum_i x_i\alpha_i$. Die übrigen Doppelmoduleigenschaften
\[
        c(y+z)=cy+cz,
        \qquad c(y\alpha)=(cy)\alpha
\]
sind trivial. Also ist $\mM$ Darstellungsmodul, der wegen der Konstruktion genau zur Darstellung $\mo\sim\mathfrak D$ gehört.

Sei nun $(y_1,\ldots,y_n)$ eine andere Basis für denselben Darstellungsmodul:
\[
        (y_1,\ldots,y_n)=(x_1,\ldots,x_n)P,
        \qquad
        (x_1,\ldots,x_n)=(y_1,\ldots,y_n)P^{-1}.
\]
Ist $b\sim B$ in bezug auf die Basis $x$, so wird
\[
        b(y_1,\ldots,y_n)=b(x_1,\ldots,x_n)P
        =(x_1,\ldots,x_n)BP
        =(y_1,\ldots,y_n)P^{-1}BP.
\]
Man nennt die Darstellungen $b\sim B$ und $b\sim P^{-1}BP$ äquivalent und rechnet sie zur selben Darstellungsklasse. Da jede reguläre Matrix $P$ eine Basis von $\mM$ wieder in eine Basis überführt, ist bewiesen:

\emph{Jeder Darstellungsmodul führt eindeutig zu einer bestimmten Darstellungsklasse.}\textsuperscript{15a)}

{\footnotesize\noindent 15) Darstellungsmoduln in bezug auf kommutative Körper finden sich zuerst bei W. Krull, \emph{Theorie und Anwendung der verallgemeinerten Abelschen Gruppen}, Sitzungsber. Heidelberger Ak. 1926, 1. Abhandl.\par
15a) Zusatz bei der Korrektur (14. April 1929). Wie B. L. van der Waerden mir mitteilt, kann man einen von der speziellen Basiswahl unabhängigen, also invarianten Zusammenhang direkt gewinnen durch Trennung der Begriffe: lineare Transformation und Matrix. Eine lineare Transformation ist ein Homomorphismus zweier Linearformenmoduln; eine Matrix ist der Ausdruck (die Darstellung) dieses Homomorphismus bei einer bestimmten Basiswahl. I. Jeder Darstellungsmodul ordnet dem Multiplikatorenbereich $\mo$ eindeutig ein homomorphes System linearer Transformationen des Moduls in sich zu. Denn nach §2 Schluß erzeugen die Linksmultiplikatoren eines Doppelmoduls $\mM$ Operatorhomomorphismen von $\mM$ in sich in bezug auf die Rechtsoperatoren (hier Multiplikation mit $K$); und zwar ist die Zuordnung eine homomorphe. II. Jedes zu $\mo$ homomorphe System linearer Transformationen eines Linearformenmoduls in sich führt zu einem Darstellungsmodul.\par}

Klar ist: Zwei operatorisomorphen Darstellungsmoduln entspricht dieselbe Darstellungsklasse. Aber auch umgekehrt: Wenn zwei Darstellungsmoduln $(x_1,\ldots,x_n)$ und $(\bar x_1,\ldots,\bar x_n)$ dieselbe Darstellung erzeugen, so gibt die Zuordnung
\[
        x_i\mapsto\bar x_i,
        \qquad
        \sum_i x_i\alpha_i\mapsto\sum_i\bar x_i\alpha_i
\]
einen Isomorphismus ab. Denn die Multiplikationsvorschrift ist durch die Matrizes $C$ vollständig bestimmt.

Spezialfall: Erzeugen zwei Basen von $\mM$ dieselbe Darstellung, so gehen sie durch einen Operatorisomorphismus von $\mM$ auf sich auseinander hervor.

Oder: Ist $C=PCP^{-1}$ für alle $C$ und festes $P$, so ist die Zuordnung
\[
        y_i\mapsto x_i,
        \qquad y_i=\sum_kx_k\pi_{ik},
\]
ein Isomorphismus von $\mM$ auf sich. Umgekehrt: Jeder Isomorphismus $y_i\mapsto x_i$ des Doppelmoduls $\mM$ auf sich wird vermittelt durch eine Matrix $P$, die mit allen $C$ vertauschbar ist.

Man kann den Darstellungsbegriff auch ausdehnen auf solche Ringe, die noch mit einem kommutativen Multiplikatorenbereich $P$ kommutativ verknüpft sind (§8). Man nimmt in diesem Fall nur solche Darstellungsringe $K$, die $P$ in ihrem Zentrum enthalten, und verlangt von der Darstellung nicht nur, daß sie ein Ringhomomorphismus $\mo\sim\mathfrak D$, sondern sogar, daß sie ein Operatorhomomorphismus sein soll: aus $a\sim A$ soll folgen
\[
        a\rho\sim A\rho\qquad(\rho\in P).
\]
Für die Darstellungsmoduln bedeutet diese Forderung die hinzukommende Rechnungsregel
\[
        am\cdot\rho=a(m\rho)=a\rho\cdot m.
\]
Der Modul und der Ring sind dann, wie man sagt, ``kommutativ mit $P$ verbunden''.

\subsection*{§ 16. Reduzible Darstellungen}
Wenn $\mU$ ein Untermodul des Darstellungsmoduls $\mM$ ist und wenn es möglich ist, eine Basis für $\mM$ zu wählen, bestehend aus einer Basis $z_1,\ldots,z_t$ von $\mU$, ergänzt durch $y_1,\ldots,y_r$, also
\[
        \mM=y_1K+\cdots+y_rK+z_1K+\cdots+z_tK,\textsuperscript{16)}
\]
so sehen die Darstellungen so aus:
\[
        C=\begin{pmatrix}R&0\\ S&T\end{pmatrix},
\]
wo die Matrizes $T$ für sich allein eine Darstellung $t$-ten Grades bilden, die durch $\mU$ erzeugt wird, und die Matrizes $R$ eine Darstellung $r$-ten Grades, die durch $\mM/\mU$ erzeugt wird.

Beweis. Es ist
\[
        (cy_1,\ldots,cy_r,cz_1,\ldots,cz_t)
        =(y_1,\ldots,y_r,z_1,\ldots,z_t)C.
\]
Da die $cz_i$, als Elemente von $\mU$, sich durch die $z_i$ allein ausdrücken, so steht in $C$ rechts oben Null. Nennt man die übrigen Abteilungen von $C$ in der obigen Anordnung $R,S,T$, so wird insbesondere
\[
        (cz_1,\ldots,cz_t)=(z_1,\ldots,z_t)T;
\]
also bilden die $T$ eine Darstellung, vermittelt durch $\mU$. Weiter wird
\[
        (cy_1,\ldots,cy_r)\equiv (y_1,\ldots,y_r)R\pmod{\mU},
\]
während die $y_i$ eine modulo $\mU$ linear-unabhängige Basis bilden. Also bilden die $R$ eine Darstellung, erzeugt durch $\mM/\mU$.

Ist umgekehrt eine ``reduzible'' Darstellung
\[
        C=\begin{pmatrix}R&0\\ S&T\end{pmatrix}
\]
gegeben, wo $R$ und $T$ quadratische Matrizes sind, so drücken sich im zugehörigen Darstellungsmodul die Produkte eines jeden $c$ mit den letzten $t$ Basiselementen $z_1,\ldots,z_t$ allein aus, d. h. $\mU=(z_1,\ldots,z_t)$ ist ein Untermodul.

{\footnotesize\noindent 16) Anders ausgedrückt: Wenn $\mU$ als $K$-Modul direkter Summand ist. Ist $K$ ein Körper, so ist diese Voraussetzung immer erfüllt.\par}

Folgerung. Sei
\[
        \mM\supset \mU_1\supset \mU_2\supset\cdots\supset \mU_e=0
\]
eine Kompositionsreihe für $\mM$, und sei jedes $\mU_i$ im vorangehenden direkter Summand als $K$-Modul, so daß man eine Basis
\[
        z_{11},\ldots,z_{1r_1},\ z_{21},\ldots,z_{2r_2},\ldots,
        z_{e1},\ldots,z_{er_e}
\]
wählen kann, so daß die $z_{ik}$ mit $i>\nu$ eine Basis für $\mU_\nu$ bilden. Dann sehen die durch $\mM$ vermittelten Darstellungen so aus:
\[
 C=\begin{pmatrix}
 R_{11}&0&\cdots&0\\
 R_{12}&R_{22}&\cdots&0\\
 \vdots&\vdots&\ddots&\vdots\\
 R_{1e}&R_{2e}&\cdots&R_{ee}
 \end{pmatrix},                                      \tag{2}
\]
und die $R_{\nu\nu}$ bilden Darstellungen, vermittelt durch die Kompositionsfaktoren $\mU_{\nu-1}/\mU_\nu$.

Die einzelnen Darstellungen $R_{\nu\nu}$ sind irreduzibel, da die Kompositionsfaktoren $\mU_{\nu-1}/\mU_\nu$ einfach sind. Setzt man voraus, daß $K$ ein Körper ist, so führt auch umgekehrt jede in der Gestalt (2) möglichst weit reduzierte Darstellung zu einer Kompositionsreihe. Nach dem Satz von Jordan-Hölder sind die Kompositionsfaktoren bis auf Operatorisomorphie eindeutig, also: die $R_{\nu\nu}$ sind bis auf äquivalente Darstellungen und bis auf die Reihenfolge eindeutig bestimmt.

Ist $\mM=\mA+\mB$ direkte Summe zweier Darstellungsmoduln $\mA=(y_1,\ldots,y_r)$, $\mB=(z_1,\ldots,z_s)$, so sieht die durch die Basis $(y_1,\ldots,y_r,z_1,\ldots,z_s)$ vermittelte Darstellung offenbar so aus
\[
        C=\begin{pmatrix}R&0\\0&S\end{pmatrix},
\]
wo $R$ die durch $\mA$, $S$ die durch $\mB$ vermittelte Darstellung des Elements $c$ bedeutet, und umgekehrt. Stellt man den Modul $\mM$ zunächst als direkte Summe von direkt unzerlegbaren Bestandteilen dar und sucht man zu diesen die Kompositionsreihen wie oben, so sieht bei passender Basiswahl die Darstellung entsprechend blockweise aus. Ist wieder $K$ ein Körper, so führt auch umgekehrt jede möglichst weit reduzierte Darstellung zu einer Darstellung des Moduls durch direkt-unzerlegbare Summanden und Kompositionsfaktoren in diesen. Nach dem Satz von Krull und Otto Schmidt sind die Klassen der direkt-unzerlegbaren Darstellungsbestandteile bis auf die Reihenfolge eindeutig bestimmt.

Ist der Modul $\mM$ vollständig reduzibel, so bestehen alle Kästen aus einer einzigen irreduziblen Darstellung, und man nennt die Darstellung vollständig reduzibel.

\subsection*{§ 17. Direkte Summenzerlegung von Darstellungsmoduln bei Ringen mit Einheitselement}
Sei $\mM$ ein $\mo$-Modul, $\mo$ Ring mit Einheit. Dann ist
\[
        \mM=\mM_e+\mM_0,
\]
wo für $\mM_e$ die Einheit Einheitsoperator, für $\mM_0$ Nulloperator ist. ($\mM_e$ und $\mM_0$ sind auch Rechts-$K$-Moduln, wenn $\mM$ es ist.)

Beweis. Jedes $m$ von $\mM$ ist darstellbar als
\[
        m=em+(m-em).
\]
Die Gesamtheit der $em$ sei $\mM_e$, die Gesamtheit der $m-em$ sei $\mM_0$. Dann sind $\mM_e$ und $\mM_0$ selbstverständlich additive Gruppen und Rechts-$K$-Moduln, falls $\mM$ es ist. Weiter ist
\[
        r em=erm,
\]
also $\mM_e$ auch $\mo$-Modul; ebenso
\[
        r(m-em)=rm-rem=rm-rm,
\]
also $\mM_0$ $\mo$-Modul. Für die $em$ ist $e$ Einheitsoperator, für die $(m-em)$ Nulloperator. Also gehört nur die Null sowohl zu $\mM_e$ wie zu $\mM_0$; die Summe $\mM_e+\mM_0$ ist direkt.

Handelt es sich um Darstellungsmoduln, so erzeugt $\mM_0$ die triviale Darstellung, wobei jedem Element die Null zugeordnet wird. Spaltet man den Summanden $\mM_0$ ab, so bleibt die durch $\mM_e$ vermittelte Darstellung, wobei dem Einheitselement die Einheitsmatrix zugeordnet ist. Auf solche Darstellungen können und wollen wir uns daher von jetzt an beschränken.

Sei
\[
        \mo=\ma_1+\cdots+\ma_s
\]
direkte Summe von zweiseitigen Idealen,
\[
        e=e_1+\cdots+e_s
\]
die Zerlegung von $e$, und $\mM$ ein $\mo$-Modul, wo $e$ Einheitsoperator ist. Dann ist
\[
        \mM=\ma_1\mM+\cdots+\ma_s\mM
\]
direkt. Offenbar ist $\mM=\mo\mM=(\ma_1\mM,\ldots,\ma_s\mM)$. Setzt man $\mb_i=\ma_1+\cdots+\widehat{\ma_i}+\cdots+\ma_s$, so ist
\[
        \mM=(\ma_i\mM,\mb_i\mM),
\]
und diese Summe ist direkt, da $e_i$ für $\ma_i\mM$ Einheitsoperator, für $\mb_i\mM$ Nulloperator ist.

Wir beschränken uns auf Grund dieses Satzes meist auf Darstellungen von zweiseitig unzerlegbaren Ringen.

\subsection*{§ 18. Modul- und Darstellungstheorie vollständig reduzibler Ringe}
Sei $\mo$ vollständig reduzibel und zweiseitig einfach, und $\mM$ endlicher $\mo$-Modul, für den das Einheitselement von $\mo$ Einheitsoperator ist. Dann ist $\mM$ vollständig reduzibel, und die einfachen Bestandteile sind den einfachen Linksidealen $\ml_i$ operatorisomorph.

Beweis. Sei
\[
        \mo=\ml_1+\cdots+\ml_n,
        \qquad
        \mM=(\mo m_1,\ldots,\mo m_k),
\]
also
\[
        \mM=(\ldots,\ml_i m_k,\ldots).                     \tag{4}
\]
Die $\ml_i m_k$, die von Null verschieden sind, sind isomorph zu $\ml_i$, denn die Zuordnung $a\mapsto am_k$ ist ein Operatorisomorphismus. Also sind die $\ml_i m_k$ einfach. Läßt man aus der Darstellung (4) diejenigen $\ml_i m_k$, die in der Summe der vorangehenden schon enthalten sind, weg, so wird die Summe direkt.

Folge. Ist außerdem $\mM$ direkt-unzerlegbar, so ist $\mM$ einfach und isomorph zu einem $\ml_i$.

Sei nun $\Delta$ der Automorphismenkörper dieses einfachen $\mM$, $\Lambda$ der von $\ml_i$; dann gilt auf Grund des Operatorisomorphismus von $\mM$ und $\ml_i$ die Ringisomorphie $\Delta\simeq\Lambda$. Nach §1 kann man $\mM$ und $\ml_i$ als Doppelmoduln mit $\Delta$ bzw. $\Lambda$ als Rechtsbereich auffassen, und zwar sind diese Doppelmoduln auch Darstellungsmoduln. Die durch sie vermittelten Darstellungen gehen durch den Ringisomorphismus $\Delta\simeq\Lambda$ auseinander hervor. Wir können uns also für das Studium dieser Darstellungsklasse auf die durch $\ml_i$ in dessen Automorphismenkörper vermittelte Darstellung beschränken.

Sei speziell für $\ml_1$ eine Basis aus Matrizeneinheiten
\[
        \ml_1=(c_{11},\ldots,c_{n1})
\]
zugrunde gelegt, und als Realisation für den Automorphismenkörper $\Lambda$ der frühere (§14) mit $K$ bezeichnete Unterkörper von $\mo$. Stellt man
\[
        c=\sum_{i,k} c_{ik}\alpha_{ik}
\]
dar, so wird
\[
        (cc_{11},\ldots,cc_{n1})=(c_{11},\ldots,c_{n1})(\alpha_{ik}).
\]
Die durch $\ml_1$ vermittelte Darstellung ist also durch die Matrix $(\alpha_{ik})$ gegeben; für $\ml_\nu=(c_{1\nu},\ldots,c_{n\nu})$ erhält man genau dieselbe Darstellung. Ist $\mo$ direkte Summe zweiseitig einfacher Ringe, so werden der Reihe nach diese Ringe isomorph, die übrigen durch Null dargestellt.

\subsection*{§ 19. Die einfachen Kompositionsfaktoren bei Moduln und Darstellungsmoduln}
Sei $\mo$ ein Ring mit Einheit, $\mc$ sein Radikal. Bei jedem einfachen Modul und Darstellungsmodul von $\mo$ wirken die Elemente aus $\mc$ annullierend; man erhält also einen Modul oder Darstellungsmodul für $\mo/\mc$.

Beweis. $\mc\mM$ ist wieder Modul, im Fall eines Darstellungsmoduls zugleich auch Rechts-$K$-Modul. Da $\mM$ einfach ist, ist entweder $\mc\mM=\mM$ oder $\mc\mM=0$. Im ersten Fall wäre $\mc^\rho\mM=\mM$ für jeden Exponenten $\rho$, im Widerspruch zur Nilpotenz von $\mc$. Also ist $\mc\mM=0$.

Daraus folgt: die einfachen Kompositionsfaktoren beliebiger Moduln und Darstellungsmoduln werden ebenfalls durch einfache Linksideale von $\mo/\mc$ erzeugt. Bei Darstellungsmoduln, die mit $P$ kommutativ verbunden sind, bedeutet das insbesondere: Reduziert man eine beliebige Darstellung mittels Kompositionsreihe wie in §16, so treten auf der Diagonale, außer Nullen, nur die durch einfache Linksideale von $\mo/\mc$ vermittelten Darstellungen auf.

Dabei braucht man für diese Aussagen keine über die bisherigen Voraussetzungen der Darstellung über $P$ hinausgehende Endlichkeitsbedingung an $\mo$. Jede Darstellung ist isomorphe Darstellung ihres absoluten Multiplikatorenbereichs; dieser ist als Urbild eines endlichen Matrizenrings selbst endlicher $P$-Modul, und da die zugelassenen Ideale $P$-Moduln sind, gelten für ihn Maximal- und Minimalbedingung.

Ist $P$ kommutativer Körper, so ist dieser absolute Multiplikatorenbereich ein hyperkomplexes System über $P$. Diese Voraussetzung ist automatisch erfüllt, sobald auf der Diagonale von Null verschiedene Darstellungen vorkommen: dann ist $P$ isomorph zu einem Unterkörper des Automorphismenkörpers eines einfachen Linksideals, und bei der Realisierung durch den Unterkörper $K$ von $\mo/\mc$ zwingt die kommutative Verbindung $P$ in das Zentrum von $K$.

\subsection*{IV. Kapitel. Darstellungen von Gruppen und hyperkomplexen Systemen}

\subsection*{§ 20. Einbeziehung der hyperkomplexen Systeme}
Wir betrachten jetzt hyperkomplexe Systeme $\mo$ in bezug auf einen kommutativen Körper $P$. Als Ideale werden verabredungsgemäß nur $P$-Moduln gerechnet. Also gelten Maximal- und Minimalbedingung für Links- und Rechtsideale, und die ganze Ringtheorie des II. Kapitels ist anwendbar.

Im folgenden soll unter Darstellung des hyperkomplexen Systems $\mo$ immer eine Darstellung im Körper $P$ verstanden sein, und zwar eine solche, für welche aus $a\sim A$ folgt
\[
        a\rho\sim A\rho\qquad(\rho\in P),
\]
also $P$-Modulhomomorphie. Für Darstellungsmoduln heißt das
\[
        a\rho\cdot m=a(m\rho).
\]
Die Ergebnisse von §19 sind also anwendbar. Alle irreduziblen Darstellungen werden durch einfache Linksideale von
\[
        \mo_0=\mo/\mc
\]
vermittelt, wo $\mc$ das Radikal ist. Es gibt also so viele inäquivalente irreduzible Darstellungen, als in der Zerlegung
\[
        \mo_0=\ma_1+\cdots+\ma_s
\]
zweiseitig einfache Summanden vorkommen.

Um die durch ein solches Linksideal vermittelte Darstellung ausdrücklich zu bestimmen, geht man zunächst von den Elementen von $\mo$ zu ihren Restklassen in $\mo_0$ über. Multiplikation mit einem Element von $P$ ist Isomorphismus für alle Ideale von $\mo_0$; also ist $P$ zu einem Unterkörper $e^{(\nu)}P$ des Automorphismenkörpers $K_\nu$ jedes einfachen Linksideals isomorph. Die durch $\ml_\nu$ über $P$ vermittelte Darstellung findet man, indem man zunächst die Darstellung über diesem Unterkörper $e^{(\nu)}P$ betrachtet und sie dann durch den Isomorphismus $e^{(\nu)}P\simeq P$ nach $P$ überträgt.

Sei
\[
        c=c_1+\cdots+c_s
\]
die Zerlegung eines Elements von $\mo_0$ nach zweiseitigen Komponenten. Für die durch $\ml_\nu$ vermittelte Darstellung ist nur die Matrix von $c_\nu$ zu betrachten, da die übrigen $c_i$ das Linksideal $\ml_\nu$ annullieren und also durch Nullmatrizes dargestellt werden. Schreibt man
\[
        c_\nu=\sum c_{ik}^{(\nu)}\alpha_{ik}^{(\nu)},
\]
so ist jedes $\alpha_{ik}^{(\nu)}$ durch die Matrix zu ersetzen, die es bei der Darstellung von $K_\nu$ über $P$ vermittelt.

Da $K_\nu$ endlichen Rang in bezug auf $P$ hat, genügt jedes Element $x$ aus $K_\nu$ einer Gleichung $f(x)=0$ mit Koeffizienten aus $e^{(\nu)}P$. Ist $P$ algebraisch-abgeschlossen, so zerfällt diese Gleichung in Linearfaktoren; da $K_\nu$ Körper ist, genügt $x$ schon einem dieser Linearfaktoren und liegt also in $e^{(\nu)}P$. Damit ist $K_\nu=e^{(\nu)}P$, und die Matrizes $(\alpha_{ik}^{(\nu)})$ selbst sind die irreduziblen Darstellungen.

Aus dem Schluß von §19 folgt damit der Satz von Burnside: Ist $P$ algebraisch-abgeschlossen und kommutativ mit $\mo$ verbunden, so enthält jede irreduzible Darstellung $n$-ten Grades in $P$ genau $n^2$ linear unabhängige Matrizes.

Hieraus folgt sofort die gewöhnliche Fassung: Ein unter Multiplikation abgeschlossenes System von Matrizes $n$-ten Grades über $P$ ist dann und nur dann irreduzibel, wenn keine nichttriviale homogene lineare Relation
\[
        \sum \alpha_{ik}x_{ik}=0
\]
zwischen den Matrixelementen aller Matrizes besteht. Fügt man nämlich alle linearen Kombinationen hinzu, so erhält man einen Ring und einen $P$-Modul und kann diesen Ring als seine eigene Darstellung auffassen.

Ebenso erhält man den verallgemeinerten Burnsideschen Satz: Unter denselben Voraussetzungen hat eine vollständig reduzible Darstellung, die in inäquivalente Bestandteile der Grade
\[
        n_1,n_2,\ldots,n_s
\]
zerfällt, den Rang
\[
        n_1^2+n_2^2+\cdots+n_s^2
\]
in bezug auf $P$.

Ist $\mo$ ein hyperkomplexes System ohne Radikal über beliebigem Körper, so ist $\mo$ vollständig reduzibel und besitzt ein Einheitselement. Nach §§17 und 18 ist jede Darstellung von $\mo$ vollständig reduzibel. Umgekehrt hat jede vollständig reduzible Darstellung eines mit $P$ kommutativ verbundenen Rings $\mo$ als absoluten Multiplikatorenbereich ein hyperkomplexes System ohne Radikal.

Reguläre Darstellung eines hyperkomplexen Systems $\mo$ ist die durch das Einheitideal $\mo$ selbst vermittelte Darstellung; bei einem Gruppenring wählt man insbesondere die Gruppenelemente als Basis. Da alle Linksideale von $\mo/\mc$ als Kompositionsfaktoren in $\mo$ vorkommen, treten alle irreduziblen Darstellungen schon auf der Diagonale der regulären Darstellung auf, wenn man diese nach §16 mittels einer Kompositionsreihe reduziert.

Eine Darstellung ist bekannt, sobald die darstellenden Matrizes für die Basiselemente $u_1,\ldots,u_h$ bekannt sind. Ist das System ein Gruppenring, so daß die $u_i$ Gruppenelemente sind, so ist jede homomorphe Matrixdarstellung der Gruppe auch eine Darstellung des Gruppenrings. Das Problem der Darstellung einer Gruppe ist somit ein Spezialfall der Darstellung hyperkomplexer Systeme.

\subsection*{§ 21. Erweiterung des Grundkörpers. Die Darstellungen des Zentrums}
Ist
\[
        \mo=a_1P+\cdots+a_hP
\]
ein hyperkomplexes System und $\Omega$ ein Erweiterungskörper des Grundkörpers $P$, so kann man bilden
\[
        \mo\Omega=a_1\Omega+\cdots+a_h\Omega
\]
mit den alten Multiplikationsregeln für die Basiselemente $a_i$.\textsuperscript{18a)} $\mo\Omega$ ist wieder hyperkomplex in bezug auf $\Omega$.

Jede Darstellung von $\mo$ führt zu einer Darstellung von $\mo\Omega$, da die Darstellung aller Elemente bekannt ist, sobald die Darstellungen der Basiselemente $a_i$ bekannt sind.\textsuperscript{19)}

{\footnotesize\noindent 18a) Genauer gesagt: man führt neue Symbole $a_i$ mit den alten Multiplikationsregeln ein; $a_1\Omega+\cdots+a_h\Omega$ enthält dann einen zu $\mo$ isomorphen Unterring, so daß man nachträglich die $a_i$ mit den alten identifizieren kann.\par
19) Es werden natürlich wieder nur solche Darstellungen betrachtet, die $P$-Modulhomomorphismen sind; aus $a\sim A$ soll folgen $a\rho\sim A\rho$ für $\rho\in P$, entsprechend für $\Omega$.\par}

Ein Ideal oder Darstellungsmodul, sowie die entsprechende Darstellung, heißen absolut-irreduzibel, wenn sie irreduzibel bleiben nach Übergang zum algebraisch abgeschlossenen Körper.

Ist $\mo\Omega$ ohne Radikal, so ist $\mo$ es auch; denn ein nilpotentes Ideal $\mc$ in $\mo$ würde zu einem Erweiterungsideal $\mc\Omega$ in $\mo\Omega$ führen. Die Umkehrung gilt nicht.

Das Zentrum von $\mo\Omega$ wird gleich dem erweiterten Zentrum $\mZ\Omega$. Ist $\mo$ vollständig reduzibel, so ist $\mZ$ direkte Summe von Körpern:
\[
        \mZ=\mZ_1+\cdots+\mZ_s.
\]
Wenn $\mo$ bei Übergang zu $\Omega$ vollständig reduzibel bleibt, so zerfällt das neue Zentrum
\[
        \mZ\Omega=\mZ_1\Omega+\cdots+\mZ_s\Omega
\]
in $\Omega$ wieder in eine direkte Summe von Körpern, die, wenn $\Omega$ algebraisch abgeschlossen ist, Rang $1$ in bezug auf $\Omega$ haben.

Für die Darstellungen des Zentrums $\mZ$ gelten die folgenden Sätze. Jede irreduzible Darstellung eines kommutativen hyperkomplexen Systems $\mZ$ im algebraisch abgeschlossenen Körper $\Omega$ ist vom ersten Grad; oder: die irreduziblen Darstellungen sind identisch mit den Homomorphismen von $\mZ$ in $\Omega$.

Beweis. Jede irreduzible Darstellung von $\mZ$ ergibt eine solche von $\mZ\Omega$, also auch eine solche von $\mZ\Omega/\mc$, wo $\mc$ das Radikal von $\mZ\Omega$ ist. Dieses ist ein kommutatives System ohne Radikal, also nach §14 Schluß direkte Summe von Körpern. Diese sind vom ersten Grad, da $\Omega$ algebraisch-abgeschlossen vorausgesetzt wurde, und erzeugen alle irreduziblen Darstellungen. Zugleich folgt: die Anzahl der verschiedenen Homomorphismen von $\mZ$ in $\Omega$ ist gleich dem Rang von $\mZ\Omega/\mc$.

Ist speziell $\mZ$ Körper über $P$, so ist $\mZ$ dann und nur dann vollständig reduzibel, wenn $\mZ$ Erweiterung erster Art von $P$ ist. Ist
\[
        \mZ=\mZ_1+\cdots+\mZ_s
\]
vollständig reduzibel, so wird bei jeder Darstellung auch jeder Körper $\mZ_i$ homomorph abgebildet, und bei einer irreduziblen Darstellung wird einer dieser Körper isomorph, die übrigen durch Null dargestellt.

Die Anwendung auf hyperkomplexe Systeme ohne Radikal ergibt: Ist $\mo\Omega$, also auch $\mo$, ohne Radikal, so werden bei jeder absolut-irreduziblen Darstellung von $\mo$ die Zentrumselemente $z$ durch Diagonalmatrizes $E\xi$ dargestellt, und $z\mapsto\xi$ ist eine Darstellung ersten Grades von $\mZ$ in $\Omega$. Die dadurch vermittelte Beziehung zwischen den absolut-irreduziblen Darstellungsklassen von $\mo$ und denen von $\mZ$ ist eineindeutig. Die Anzahl dieser Darstellungsklassen ist also gleich dem Rang des Zentrums.\textsuperscript{20)}

{\footnotesize\noindent 20) Entsprechende Sätze gelten nicht nur für Darstellungen im algebraisch-abgeschlossenen Körper $\Omega$, sondern auch für Darstellungen in den einzelnen Automorphismenkörpern.\par}

Der Zerfällung in zweiseitig-unzerlegbare Ideale
\[
        \mo\Omega=\ma_1+\cdots+\ma_s
\]
entspricht eineindeutig eine Zerfällung des Zentrums in Körper
\[
        \mZ\Omega=\mZ_1+\cdots+\mZ_s.
\]
Bei einer irreduziblen Darstellung von $\mo$ wird ein $\ma_i$ isomorph, die übrigen durch Null dargestellt, und die Darstellungsklasse ist durch $\ma_i$ eindeutig bestimmt. $\ma_i$ wird durch einen vollen Matrizenring dargestellt, dessen Zentrum aus Diagonalmatrizen $E\xi$ besteht; die Zuordnung $z\mapsto\xi$ ist der zugehörige Homomorphismus des Zentrums.

Für die Frage, wann einem $\mo$ ohne Radikal ein ebensolches $\mo\Omega$ entspricht, folgt noch: Hat $\mZ$ eine Komponente $\mZ_i$, die Erweiterung zweiter Art von $P$ ist, so besitzt $\mo\Omega$ sicher ein Radikal.\textsuperscript{21)} Denn $\mZ_i\Omega$ besitzt in diesem Falle schon eines.

{\footnotesize\noindent 21) Hat $\mZ$ nur Komponenten von erster Art, so wird $\mo\Omega$ ohne Radikal. Für Charakteristik Null hat schon I. Schur den äquivalenten Satz bewiesen, daß irreduzible Darstellungen in $P$ bei jeder Erweiterung des Koeffizientenbereichs vollständig reduzibel bleiben.\par}

\subsection*{§ 22. Anwendung auf Abelsche Gruppen}
Sei
\[
        G=G_1\times G_2\times\cdots\times G_r
\]
eine endliche Abelsche Gruppe, zerlegt in zyklische Gruppen $G_i$ der Ordnungen $h_i$. Elemente:
\[
        a=a_1^{\alpha_1}\cdots a_r^{\alpha_r}.
\]
Es sei $P$ ein Körper, dessen Charakteristik nicht in der Gruppenordnung $h=h_1h_2\cdots h_r$ aufgeht.

Der Gruppenring $\mZ$ besteht aus allen Summen
\[
        \sum A_{\alpha_1\ldots\alpha_r}a_1^{\alpha_1}\cdots a_r^{\alpha_r},
        \qquad A_{\alpha_1\ldots\alpha_r}\in P,
\]
und ist homomorphes Abbild des Polynombereichs $P[x_1,\ldots,x_r]$ vermöge $x_i\mapsto a_i$. Also ist
\[
        \mZ\simeq P[x_1,\ldots,x_r]/(x_1^{h_1}-e,\ldots,x_r^{h_r}-e).
\]
Im Erweiterungskörper $\Omega$ zerfallen die $x_i^{h_i}-e$ in lauter verschiedene Faktoren. Das definierende Ideal wird Produkt oder Durchschnitt lauter verschiedener Primideale
\[
        (x_1-\xi_1^{(\alpha_1)},\ldots,x_r-\xi_r^{(\alpha_r)}),
\]
mit je nur einer Nullstelle. Der Durchschnittsbildung entspricht eine direkte Summenzerlegung
\[
        \mZ\Omega=\mZ_1+\cdots+\mZ_h,
\]
wo jeweils $\mZ_i\sim\Omega$ wird, also die Darstellung
\[
        a_i\mapsto \xi_i^{(\alpha_i)},
        \qquad\hbox{oder}\qquad a\mapsto \chi^{(\alpha)}(a)
\]
vermittelt. Diese Darstellungen sind die Charaktere. Ihre Anzahl ist $h_1\cdots h_r=h$.

\subsection*{§ 23. Determinante eines hyperkomplexen Systems}
Sei
\[
        \mo=a_1P+\cdots+a_hP
\]
ein hyperkomplexes System. Ich adjungiere zu $P$ die Unbestimmten $x_1,\ldots,x_h$ und bilde
\[
        \mo^*=a_1P(x)+\cdots+a_hP(x).
\]
Die $x_i$ sollen mit den $a_i$ vertauschbar sein; dadurch sind die Rechnungsregeln in $\mo^*$ bestimmt.

In $\mo^*$ liegt das allgemeine Element von $\mo$,
\[
        w=a_1x_1+\cdots+a_hx_h.
\]
Ist in einer Darstellung $a_i\mapsto A_i$, so ist dem $w$ zugeordnet
\[
        W=A_1x_1+\cdots+A_hx_h.
\]
$W$ heißt die zur Darstellung gehörige Systemmatrix; speziell im Gruppenring die Gruppenmatrix. Handelt es sich um die reguläre Darstellung, so hat man die reguläre Systemmatrix.

Die Elemente von $W$ sind Linearformen der $x_i$. Die Systemdeterminante $|W|$ ist also vom Grad $n$, wenn es sich um Darstellungen $n$-ten Grades handelt. Insbesondere ist die reguläre Systemdeterminante vom Grad $h$.

Die Systemdeterminante ändert sich nicht bei Übergang zu äquivalenten Darstellungen, da $|PWP^{-1}|=|P|\,|W|\,|P^{-1}|=|W|$. Bei Übergang von $(a_1,\ldots,a_h)$ zu einer neuen Basis $(b_1,\ldots,b_h)$ und von $w=\sum a_ix_i$ zu $w=\sum b_jy_j$ findet man die neue Determinante aus der alten durch eine reguläre Substitution
\[
        x_i=\sum_j\rho_{ij}y_j.
\]

Liegt eine Kompositionsreihe des Darstellungsmoduls vor, so sieht bei passender Basiswahl die Matrix $W$ so aus:
\[
        W=\begin{pmatrix}
        W_1&0&\cdots&0\\
        *&W_2&\cdots&0\\
        \vdots&\vdots&\ddots&\vdots\\
        *&*&\cdots&W_s
        \end{pmatrix}.
\]
Unter den $|W_i|$ einer beliebigen Darstellung kommen keine anderen vor als in der regulären Darstellung von $\mo$ oder sogar von $\mo/\mc$, wo $\mc$ das Radikal ist.

Ist $P$ algebraisch abgeschlossen, so ist die zu einer irreduziblen Darstellung gehörige Determinante $|W_i|$ eine Primfunktion in den $x_i$, und zu inäquivalenten Darstellungen gehören verschiedene Primfaktoren.

Beweis. Wir können uns auf den Ring ohne Radikal $\mo/\mc$ beschränken. In ihm führen wir als Basis die Matrizeneinheiten $c_{ik}^{(\nu)}$ ein; das allgemeine Element lautet
\[
        w=\sum_{\nu,i,k}c_{ik}^{(\nu)}x_{ik}^{(\nu)}.
\]
Die Matrizes der irreduziblen Darstellungen lauten
\[
        W_\nu=(x_{ik}^{(\nu)}).
\]
Die Funktionen $|W_\nu|=|x_{ik}^{(\nu)}|$ sind irreduzibel und voneinander verschieden.

Um die $|W_i|$ zu berechnen, kann man die reguläre Systemmatrix von $\mo/\mc$ in ihre Primfaktoren zerlegen. Jeder Primfaktor kommt so oft vor, wie der Grad der irreduziblen Darstellung angibt. Man kann auch die reguläre Systemmatrix von $\mo$ zugrunde legen, erhält dann aber jeden irreduziblen Faktor öfter; bei der regulären Darstellung als Rechtsideal entsteht die antistrophe Matrix bei Frobenius, welche dieselben irreduziblen Faktoren enthält, aber möglicherweise mit anderen Exponenten.

Die Systemdeterminante eines kommutativen Systems zerfällt in Linearfaktoren, da alle irreduziblen Darstellungen vom ersten Grad sind. Diese Linearfaktoren sind selbst die irreduziblen Darstellungen; sie ergeben also bei Abelschen Gruppen die Charaktere. Diese Tatsache war Dedekinds Ausgangspunkt bei der Untersuchung der Gruppendeterminante nichtabelscher Gruppen.

\subsection*{§ 24. Spuren und Charaktere}
Ist in einer Darstellung $\mo\sim\mathfrak D$ eines hyperkomplexen Systems $\mo$ dem Element $a$ die Matrix $A$ zugeordnet, so setzt man
\[
        \operatorname{Sp}_D(a)=\operatorname{Sp} A.
\]
Spuren sind lineare Funktionen:
\[
        \operatorname{Sp}(c+d)=\operatorname{Sp}(c)+\operatorname{Sp}(d),
        \qquad
        \operatorname{Sp}(c\alpha)=\alpha\operatorname{Sp}(c).
\]
Äquivalente Darstellungen haben dieselben Spuren. Die Spur in einer reduziblen Darstellung ist Summe der Spuren in den durch die Kompositionsfaktoren vermittelten Darstellungen; denn
\[
        \operatorname{Sp}\begin{pmatrix}
        A_{11}&0&\cdots&0\\
        *&A_{22}&\cdots&0\\
        \vdots&\vdots&\ddots&\vdots\\
        *&*&\cdots&A_{ss}
        \end{pmatrix}
        =\operatorname{Sp}(A_{11})+\operatorname{Sp}(A_{22})+\cdots+\operatorname{Sp}(A_{ss}).
\]
Hauptspur heißt die Spur bei der regulären Darstellung. Reduzierte Spur heißt die Summe der Spuren bei den verschiedenen irreduziblen Darstellungen.

Ist $c$ Element des maximalen nilpotenten Ideals $\mc$, so ist $\operatorname{Sp}(c)=0$ für jede Darstellung. Es genügt, die Darstellungen durch die einfachen Linksideale von $\mo/\mc$ zu untersuchen; $c$ annulliert aber alle Elemente von $\mo/\mc$, also ist jedem $c$ aus $\mc$ die Nullmatrix zugeordnet.

Ist $\mo$ zweiseitig einfach und $P$ algebraisch-abgeschlossen, also Matrizenring $\mo=\sum c_{ik}P$, so ist bei der irreduziblen Darstellung von $\mo$ dem Element $a=\sum c_{ik}\alpha_{ik}$ die Matrix $(\alpha_{ik})$ zugeordnet. Also
\[
        \operatorname{Sp}_{\rm red}(a)=\sum_i\alpha_{ii},
        \qquad
        \operatorname{Sp}_{\rm pr}(a)=n\,\operatorname{Sp}_{\rm red}(a)=n\sum_i\alpha_{ii}.
\]
Insbesondere
\[
        \operatorname{Sp}_{\rm red}(c_{ik})=0\ (i\ne k),\quad
        \operatorname{Sp}_{\rm red}(c_{ii})=e,
\]
und
\[
        \operatorname{Sp}_{\rm pr}(c_{ik})=0\ (i\ne k),\quad
        \operatorname{Sp}_{\rm pr}(c_{ii})=n e.
\]

Geht man von $P$ zum algebraisch abgeschlossenen Körper $\Omega$ über, so ändert sich die Hauptspur nicht, da sie aus derselben Basis berechnet werden kann. Das gilt nicht für die reduzierte Spur.

Die Spuren der Elemente $a$ des Systems ohne Radikal in den absolut-irreduziblen Darstellungen nennt man Charaktere und bezeichnet sie mit $\chi(a)$ oder, wenn angegeben werden soll, welche Darstellung gemeint ist, mit $\chi^{(\nu)}(a)$.

Bei einer irreduziblen Darstellung vom Grad $n_\nu$ werden die Zentrumselemente nach §21 durch Diagonalmatrizes $E\theta^{(\nu)}$ dargestellt, wo $z\mapsto\theta^{(\nu)}(z)$ eine Darstellung ersten Grades des Zentrums ist. Daraus folgt für die Spur der Wert $n_\nu\theta^{(\nu)}(z)$. Also sind die Homomorphismen des Zentrums mit den Charakteren durch
\[
        \chi^{(\nu)}(z)=n_\nu\theta^{(\nu)}(z)                 \tag{1}
\]
verknüpft. Im kommutativen Fall ist $n_\nu=1$, und die Charaktere selbst geben die Homomorphismen. Hat der Körper $\Omega$ die Charakteristik Null, so kann man (1) durch $n_\nu$ dividieren:
\[
        \theta^{(\nu)}(z)=\frac{\chi^{(\nu)}(z)}{n_\nu}.
\]
Die Homomorphismuseigenschaft drückt sich durch
\[
        \frac{\chi^{(\nu)}(z+z')}{n_\nu}
        =\frac{\chi^{(\nu)}(z)}{n_\nu}+\frac{\chi^{(\nu)}(z')}{n_\nu},
\]
\[
        \frac{\chi^{(\nu)}(zz')}{n_\nu}
        =\frac{\chi^{(\nu)}(z)}{n_\nu}\frac{\chi^{(\nu)}(z')}{n_\nu}
\]
aus.

Eine Darstellungsklasse ist durch die Spuren der Matrizes allein schon eindeutig festgelegt. Beweis: Der Darstellungsmodul $R$ ist bekannt, wenn man weiß, wie oft jeder irreduzible Darstellungsmodul $\mM_\nu$ als direkter Summand vorkommt. Ist diese Anzahl $\mu_\nu$, so ist die Spur von $e^{(\nu)}$ gleich $\mu_\nu n_\nu$, also
\[
        \mu_\nu=\frac{\operatorname{Sp}(e^{(\nu)})}{n_\nu}.
\]

\subsection*{§ 25. Diskriminanten}
Sei
\[
        \mo=a_1P+\cdots+a_hP
\]
ein hyperkomplexes System. Diskriminantenmatrix heißt die Matrix, deren Elemente die Hauptspuren $\operatorname{Sp}(a_i a_j)$ sind. Reduzierte Diskriminantenmatrix heißt dasselbe bei reduzierten Spuren, gebildet im algebraisch-abgeschlossenen Körper. Die Determinante der Matrix heißt Diskriminante beziehungsweise reduzierte Diskriminante.

Die Diskriminante bleibt dieselbe bei Erweiterung des Grundkörpers. Bei Übergang zu einer anderen Basis multipliziert sie sich mit dem Quadrat der Transformationsdeterminante. Ist
\[
        (b_1,\ldots,b_h)=(a_1,\ldots,a_h)P,
\]
so gilt
\[
        (\operatorname{Sp}(a_i a_j))=(\operatorname{Sp}(a_i b_j))P,             \tag{1}
\]
und ebenso, wenn $P^t$ die gespiegelte Matrix ist,
\[
        (\operatorname{Sp}(a_i b_j))=P^t(\operatorname{Sp}(b_i b_j)).           \tag{1a}
\]
Also
\[
        (\operatorname{Sp}(a_i a_j))=P^t(\operatorname{Sp}(b_i b_j))P.
\]
Der Übergang zu Determinanten liefert die Behauptung. Die Determinante ist also nur bis auf eine Quadratzahl aus $P$ bestimmt; ihr Verschwinden ist eine invariante Eigenschaft. Aus (1) folgt noch, daß man die Diskriminante bis auf einen Faktor $\ne0$ auch aus zwei verschiedenen Basen als $|\operatorname{Sp}(a_i b_j)|$ bestimmen kann.

Die Diskriminante verschwindet, wenn $\mo$ ein nilpotentes Ideal $\mc$ besitzt. Als Basiselemente wähle man
\[
        c_1,\ldots,c_r,d_1,\ldots,d_s,
\]
wo $c_1,\ldots,c_r$ eine Basis von $\mc$ bilden. Dann hat die Diskriminantenmatrix einen Block mit verschwindenden Spuren,
\[
        \begin{pmatrix}
        \operatorname{Sp}(c_i c_j)&\operatorname{Sp}(c_i d_j)\\
        \operatorname{Sp}(d_i c_j)&\operatorname{Sp}(d_i d_j)
        \end{pmatrix},
\]
und der obere Teil wird Null; das gilt auch für die reduzierte Diskriminante. Die Diskriminante verschwindet daher auch dann, wenn erst nach Übergang zum algebraisch-abgeschlossenen Körper ein nilpotentes Ideal auftritt.

Sei
\[
        \mo=\ma_1+\cdots+\ma_s
\]
direkt zweiseitig, und seien $M,M_1,\ldots,M_s$ die Diskriminantenmatrizes der Ringe $\mo,\ma_1,\ldots,\ma_s$. Dann wird
\[
        M=\begin{pmatrix}
        M_1&0&\cdots&0\\
        0&M_2&\cdots&0\\
        \vdots&\vdots&\ddots&\vdots\\
        0&0&\cdots&M_s
        \end{pmatrix},
\]
also $|M|=|M_1|\cdots|M_s|$. Wählt man eine Basis für $\mo$, die sich aus den Basen der $\ma_i$ zusammensetzt, so ist für $a\in\ma_i$ die Hauptspur in $\mo$ gleich der Hauptspur in $\ma_i$; bei verschiedenen Komponenten verschwinden die Mischprodukte.

Für einen Matrizenring $\sum c_{ik}P$ erhält man, wenn die Matrizeneinheiten einmal nach ersten und einmal nach zweiten Indizes geordnet werden,
\[
        D_{\rm red}=e,
        \qquad
        D=n^{n^2}e.
\]
Zerfällt $\mo$ im algebraisch-abgeschlossenen Erweiterungskörper in Matrizenringe der Grade $n_1,\ldots,n_s$, so wird
\[
        D=n_1^{n_1^2}n_2^{n_2^2}\cdots n_s^{n_s^2}e,
        \qquad
        D_{\rm red}=e.
\]
Die reduzierte Diskriminante ist demnach bei Systemen ohne Radikal immer $\ne0$, die andere nur dann, wenn kein $n_i$ durch die Charakteristik des Körpers teilbar ist. Mithin: Das Nichtverschwinden der reduzierten Diskriminante ist notwendig und hinreichend für Systeme ohne Radikal nach algebraischem Abschluß des Körpers $P$. Im Fall der Charakteristik Null ist das Nichtverschwinden der Diskriminante notwendig und hinreichend für das Nichtauftreten eines Radikals nach algebraischem Abschluß.\textsuperscript{22)}

{\footnotesize\noindent 22) Vgl. für kommutative Systeme E. Noether, \emph{Diskriminantensatz für Ordnungen ...}, J. f. M. 157 (1927), S. 82--104, §§4--6. Die dortige Beweismethode ist die gleiche, aber im einzelnen komplizierter; der Übergang von einer Basis zu einer andern ist durch den hier gegebenen zu ersetzen.\par}

\subsection*{§ 26. Einordnung des Gruppenrings}
Sind $a_1,\ldots,a_h$ die Elemente einer endlichen Gruppe und bildet man den Gruppenring in einem Körper, dessen Charakteristik nicht Teiler von $h$ ist, so ist die Diskriminante $D\ne0$, also der Gruppenring ein Ring ohne Radikal.

Beweis. Wir zeigen zuerst, wobei $\operatorname{Sp}$ fortan Hauptspur heißt:
\[
        \operatorname{Sp}(e)=he,
        \qquad
        \operatorname{Sp}(a_i)=0\quad(a_i\ne e).
\]
In der regulären Darstellung ist nämlich
\[
        e\mapsto\begin{pmatrix}e&0&\cdots&0\\0&e&\cdots&0\\\vdots&\vdots&\ddots&\vdots\\0&0&\cdots&e\end{pmatrix},
        \qquad \operatorname{Sp}(e)=he.
\]
Für $a_i\ne e$ ist $a_i a_1,\ldots,a_i a_h$ eine Permutation ohne Fixpunkt der Basiselemente; die entsprechende Matrix besitzt in der Hauptdiagonale nur Nullen, also $\operatorname{Sp}(a_i)=0$.

Wir nehmen nun die Basen $a_1,\ldots,a_h$ und $a_1^{-1},\ldots,a_h^{-1}$. Die Matrix $(\operatorname{Sp}(a_i a_j^{-1}))$ lautet diagonal mit $he$ in der Hauptdiagonale, also
\[
        D=h^h e\ne0.
\]

Als Klasse eines Gruppenelements bezeichnet man die im Gruppenring gebildete Summe
\[
        K_a=\sum_s s^{-1}as,                         \tag{1}
\]
wobei nur über die verschiedenen $s^{-1}as$ summiert wird. Die Klassen $K_a$ sind Zentrumselemente, da sie mit allen Gruppenelementen vertauschbar sind. Sie erzeugen das Zentrum; denn wenn ein Element $\sum_i\alpha_i a_i$ des Gruppenrings mit allen $s$ vertauschbar ist, so ist
\[
        \sum_i\alpha_i a_i=s^{-1}\biggl(\sum_i\alpha_i a_i\biggr)s
        =\sum_i\alpha_i(s^{-1}a_i s),
\]
also sind die Koeffizienten auf jeder Konjugiertenklasse gleich.

Folglich ist der Rang des Zentrums gleich der Anzahl der Klassen konjugierter Gruppenelemente, mithin auch die Anzahl der absolut-irreduziblen Darstellungen gleich dieser Klassenzahl. Die Matrizes $A$ und $S^{-1}AS$ haben dieselbe Spur; also haben die Gruppenelemente $a$ und $s^{-1}as$ dieselbe Spur. Bildet man auf beiden Seiten von (1) die Spur, so folgt
\[
        \chi(K_a)=h_a\chi(a),
\]
wo $h_a$ die Anzahl der Elemente der Klasse $K_a$ ist. In Worten: Der Charakter eines Gruppenelements ist gleich dem Charakter der Klasse, dividiert durch die Anzahl der Elemente der Klasse.

\begin{center}
(Eingegangen am 12. August 1928.)
\end{center}

\end{document}
